\section*{Diagnostics and excluded shortcuts}
On a smooth compact curve, every nontrivial semivaluation has the form
\(t\operatorname{ord}_p\), \(0<t\leq+\infty\)
\cite[Example~2.3.3]{MP24}.
Positive rational \(t_j\leq t\) tending to \(t\), with \(t_j\to+\infty\)
when \(t=+\infty\), give the claimed approximation. At the trivial
valuation one uses the trivial valuation itself; omitting it from the
divisorial set would make the assertion false there.

For a local surface diagnostic, take a monomial semivaluation with
weights \((\sqrt2,+\infty)\) on coordinates \((x,y)\). For
\(\mathfrak a=(x^3,y^2)\) and \(\mathfrak b=(y)\), its values are
\(3\sqrt2\) and \(+\infty\). Rational weights
\(r_j\uparrow\sqrt2\) and \(M_j\to+\infty\) give
\[
 q_j(\mathfrak a)=\min\{3r_j,2M_j\}\longrightarrow3\sqrt2,\qquad
 q_j(\mathfrak b)=M_j\longrightarrow+\infty.
\]
The same coordinatewise inequality gives domination on every
convergent power series, not merely on these two test ideals.
The global proof uses principalization and the annihilator in Step 2
to reproduce this local reasoning for all global coherent ideals.

There is a genuine constraint in the tempting test-configuration
argument. If a normalized simplex has equation
\(\sum_i b_i w_i=1\), then \(w'_i\leq w_i\) for all \(i\), with the
same normalization, forces \(w'=w\). For example, when
\(b_1=b_2=1\), an irrational point
\((\sqrt2/2,1-\sqrt2/2)\) admits no distinct coordinatewise smaller
normalized rational point. This does not disprove dominated density
on \(X^\beth\): domination there tests ideals on \(X\), not arbitrary
vertical divisors of one fixed test configuration.
The proof deliberately works on log resolutions of \(X\), where
monomial weights lie in a cone without that normalization constraint.

Ordinary divisorial density alone also does not imply the conclusion
for an arbitrary ordered space. It is used only for the closed-image
surjectivity in Step 1. Step 2 provides the additional order control.
No increasing divisorial net, local description of every global
semivaluation by a valuation at a closed point, or unique lifting of
semivaluations is asserted.

